\documentclass{article}
\usepackage{amsmath}
\usepackage{siunitx}
\usepackage{physics}
\usepackage{geometry}
\geometry{a4paper, margin=1in}

\title{Shear Amplitude Calculation in Abaqus Simulation}
\author{}
\date{}

\begin{document}

\maketitle

\section*{Parameter Definitions}

\begin{itemize}
    \item $\mu$: Friction coefficient, unitless ($\mu = 0.70$)
    \item $\sigma_n$: Normal stress, in MPa ($\sigma_n = \SI{10}{\mega\pascal}$)
    \item $Y$: Yield strength of PMMA, in MPa ($Y = \SI{3000}{\mega\pascal}$)
    \item $A$: Contact area between spring and plate, $A = \texttt{SPR\_W} \times \texttt{SPR\_D}$, in mm$^2$
    \item $h$: Height of the spring ($\texttt{SPR\_H}$), in mm
    \item $D$: Depth of contact interface ($\texttt{DEPTH}$), in mm
    \item $L$: Resistance area length, I use \SI{200}{\milli\meter}
\end{itemize}

\section*{Step 1: Frictional Resistance Force}

\begin{equation}
F_{\text{resist}} = 2 \mu \cdot \sigma_n \cdot D \cdot L
\end{equation}

\textbf{Unit Check:}
\begin{align*}
[\text{MPa}] \cdot [\text{mm}] \cdot [\text{mm}] &= \frac{\si{\newton}}{\si{\milli\meter\squared}} \cdot \si{\milli\meter} \cdot \si{\milli\meter} \\
&= \si{\newton}
\end{align*}

So, $F_{\text{resist}}$ has unit of \si{\newton}.

\section*{Step 2: Shear Amplitude}

The shear amplitude is computed from:

\begin{equation}
\delta = \frac{F_{\text{resist}} \cdot h}{A \cdot Y}
\end{equation}

\textbf{Unit Check:}
\begin{align*}
\frac{[\si{\newton}] \cdot [\si{\milli\meter}]}{[\si{\milli\meter\squared}] \cdot [\si{\mega\pascal}]} 
&= \frac{\si{\newton} \cdot \si{\milli\meter}}{\si{\milli\meter\squared} \cdot \frac{\si{\newton}}{\si{\milli\meter\squared}}} \\
&= \si{\milli\meter}
\end{align*}

Thus, $\delta$ (shear amplitude) has unit of \si{\milli\meter}, which is physically consistent as a displacement.

\section*{Alternative Expression (Code Form)}

In Python-style code:

\begin{verbatim}
RESISTANCE = 2 * FRICTION_COEFFICIENT * DEPTH * 220 (mm) * NORMAL_STRESS
SHEAR_AMPLITUDE = RESISTANCE * SPR_H / ( SPR_W * SPR_D * Y_PMMA )
\end{verbatim}

Or, more explicitly:

\begin{verbatim}
CONTACT_AREA = SPR_W * SPR_D
SPRING_STIFFNESS = Y_PMMA * CONTACT_AREA / SPR_H
SHEAR_AMPLITUDE = RESISTANCE / SPRING_STIFFNESS
\end{verbatim}

\end{document}